\section{Local observations of the electrical network}
\label{sec:locality}

The geometric convergence of Proposition~\ref{prop:geometry} does not
determine the conductance of a set of edges.  We shall obtain local
electrical information by retaining numerical functions of finite pieces
of the map in the conditional-law construction of the companions.  This
section specifies the functions that can be retained and the changes of
surface law that preserve their almost sure properties.

We use the protected charts of the canonical companion
\cite[Section~5]{C} and, separately, those of the metric-measure companion
\cite[Sections~5 and~7]{P}.  A chart instruction consists of finitely
many used time intervals, each strictly inside an extended raw-word
block; rational relative time cells for ports, confinement, suppliers,
and the stars to be retained; exact differences of the two contour
heights at used starts, recorded on a forest on the block indices in each
coordinate; and finite instructions for gluing the retained pieces.
These last instructions give rational height bands for cross-block
contacts, exact lattice minima on retained subintervals, and types for
externally unresolved flexible orders in the used portions.  Overlapping
extended blocks are merged before the instruction is formed; the
intervening unused time intervals are its gaps.  Durations and height differences are
parameters, whereas the combinatorial instruction belongs to a countable
list.

Write \(\mathcal R_m^\alpha\) for the resulting finite incidence record
at \(m\) letters per time unit, for instruction \(\alpha\).  In the
canonical construction this record includes the oriented flag gluings
and every incidence occurrence and cyclic order at each protected primal
and dual vertex.  It therefore determines both edge networks there,
including loops and repeated edges.  The protected-atlas statement of the
metric-measure companion gives the primal incidences used by its tests.
We use the canonical record whenever complete dual incidences are needed.

The continuous input \(I^\alpha\) consists of the directed local
traversals, their quantum durations and labelled frontier lengths, the
cut marks, and their finite relative traversal order.  Here the
traversal and frontier data include, in word coordinates, the rescaled
increment paths on the full extended blocks, their durations and
used-start offsets, and the scaled forest differences.  A common clock
origin and a common height origin are removed.  Let \(\mathcal C\)
denote the complete continuum variable: the limiting contours, their
curve-decorated surface, the fields and cut marks when used, and the
chosen coordinate data.  In the canonical protocol, any extra embedding
or imaginary-field lift is sampled from its conditional law given the
curve-decorated continuum surface, independently of the discrete outputs
conditional on that surface; see \cite[Theorem~5.9]{C}.  In the
metric-measure protocol, additional coordinates not reconstructed from
the complete contour-and-cut data are sampled from the canonical
conditional law given those data, independently of the entire actual and
auxiliary-reference observation array conditional on those data; see
\cite[Theorem~7.10]{P}.
The chosen coordinate data in \(\mathcal C\) consist only of data
reconstructed from the appropriate prescribed continuum variable or
attached by its stated output-independent protocol.  The canonical
uniformization comparison and its orientation sign may be retained in the
joint coupling.  If their conditional law depends on the observation
array beyond the prescribed continuum input, they are kept outside the
kernel-conditioning sigma-field \(\sigma(\mathcal C)\).  The marginal
almost sure statements are transferred before that comparison and sign
are used pathwise.

For each matching comparison on a protected real-field buffer \(U\), fix
a Gaussian reference with its additive height retained: either \(G_R|_U\),
where a neighborhood of \(\overline U\) lies in \(B_R\) and \(G_R\) is a
whole-plane Gaussian free field (GFF) with zero circle average on
\(\partial B_R\), or the restriction of a zero-boundary GFF from a fixed
containing disk.  These covariance kernels have logarithmic singularity
\(\log|z-w|^{-1}\).  These choices, fixed smooth shifts on a neighborhood
of \(\overline U\), and cone restrictions on buffers strictly inside their
normalization circles and away from their marked centers are locally mutually
absolutely continuous on strict
buffers; they are not identical laws
\cite[Proposition~5.12]{C} \cite[Lemmas~3.3 and~9.2]{P}.
We call a chart carrying one such fixed reference an \emph{ordinary field
chart}.  Before the area clock, it is accompanied by an independent
whole-plane imaginary field modulo \(2\pi\chi\mathbb Z\),
\(\chi=2/\gamma-\gamma/2\), with uniform phase in \([0,2\pi\chi)\), and
its unparameterized directed curve, as in Section~3 of~\cite{C}, and by
the conditional local Poisson area-mark law \(\Pi_h\) of intensity
\(\mu_h(\dd z)\,\dd u\) on \(U\times(0,\infty)\), as in Section~7
of~\cite{P}.  These Poisson marks are separate from the auxiliary sphere
area marks and the three canonical flag samples.

A numerical observation at \(\alpha\) has the form
\begin{equation}
 Y_m^\alpha
   =F_m^\alpha(\mathcal R_m^\alpha,\theta_m^\alpha)\in E_\alpha ,
 \label{eq:finite-readout}
\end{equation}
where \(\theta_m^\alpha\) comprises the exact parameters of the
instruction, \(E_\alpha\) is a compact metric space, and \(F_m^\alpha\)
is measurable.  The same function is used in the actual word and in its
reference experiment.  Its deterministic dependence on \(m\) may include
any prescribed positive unit.  A nonnegative extended-real quantity is
compactified by \(x\mapsto x/(1+x)\), with \(\infty\mapsto1\).
The observation uses only the record and the intrinsic instruction.  In
particular, it uses no absolute clock origin, omitted-gap duration, or
edge whose incidences have not been retained.  It may be given any fixed
value when the finite instruction is not realized.

Spatial meaning is assigned with margins.  For \(K\Subset U\) in a
reference coordinate, choose
\begin{equation}
 K\subset\operatorname{int}K_*
   \Subset U_0\Subset U_1\Subset U .
 \label{eq:local-protections}
\end{equation}
The used interiors cover the full preimage of \(K_*\), their images have
closures in \(U_0\), and the extended blocks have closures in \(U_1\).
All stars used by the observation are included.  Work on the common
probability-one Brownian generic event of the reconstruction results.
We say the instruction is \emph{activated} when the containments, local
band-separation tests, and local flexible-order instructions hold for its
limiting input.  Activation is an event of the limiting local traversals,
frontier lengths, marks, and finite order.  The global generic event and
the enclosing supplier windows are used to prove eventual correctness,
not to select the instruction.  Protected reconstruction then makes
every tested incidence exact for all sufficiently large discrete indices.
Finite inner and outer cell covers between the sets in
\eqref{eq:local-protections} describe ports and confinement.  Their
strict margins are fixed before the discrete limit.

\begin{proposition}[Extension to numerical incidence observations]
\label{prop:readout-extension}
Fix either companion's chart protocol, a countable collection of its
intrinsic instructions, and observations \eqref{eq:finite-readout}.
Assume the following for every instruction that is used.
\begin{enumerate}[label=(\roman*)]
\item Activation is selected from the limiting local data above and
supplies positive limiting durations, extensions and gaps, and the
complete preimages and stars needed by the observation, with strict
inner and outer margins.
\item Descriptions of the same observation are restrictions of one
common finite union incidence record on activation.  Charts in
disjoint protected buffers are sampled as independent complete
reference packages before a resampling weight is applied.
\item Actual and reference observations use the same exact parameters
and deterministic units.  In an indexed extraction, each fixed
retained effective time unit tends to infinity in lattice units.
Varying labels are re-expressed with a fixed normalized ordinary
reference input law and positive duration and shift densities in that
unit convention, as in the companion's indexed construction.
\end{enumerate}
From every sequence of scales tending to infinity one can pass to a further joint
representation extraction, retaining these observations and
\(\mathcal C\), with the following properties.

For each instruction there is a conditional kernel
\(K^\alpha(I^\alpha,\dd y)\).  Writing \(A^\alpha\) for its activation
event, in the bilateral word and in every strict interior chart of the
finite empty-word bridge one has, for bounded measurable \(f\),
\begin{equation}
 \ind_{A^\alpha}\EE[f(Y^\alpha)\mid\mathcal C]
 =\ind_{A^\alpha}\int f(y)K^\alpha(I^\alpha,\dd y).
 \label{eq:activated-kernel}
\end{equation}
On joint activation, disjoint protected buffers have
product conditional kernels, including countable arrays in each buffer.
Descriptions of the same observation agree through their union record.
No claim is made about a default value on nonactivation.  The canonical
protocol's augmentation over all feasible finite traversal orders also
applies to the enlarged arrays, retaining the actual-order component
and conditional products in disjoint raw-buffer packages.
These conclusions persist under a fixed
locally absolutely continuous change of the continuum field law and
under the indexed changes of units described in the proof.

On activation, the kernel of a spatial observation is measurable from
the real field in \(U\), the imaginary field in a fixed larger buffer,
the local area marks, and the finite relative order of the traversals
used by the observation.  This assertion selects from the jointly
extracted limiting array; it does not condition a finite word on success
of the numerical observation.  The normalized reference law and its
matching kernels are fixed within each indexed extraction.  The two
companion protocols furnish these properties separately, without an
identification of their kernels.
\end{proposition}

\begin{proof}
We first explain why the reference-word calculation permits the added
coordinate.  This extends the proof of
\cite[Lemma~5.8 and Theorem~5.9]{C}; the fixed-chart theorem
\cite[Theorem~7.10]{P} already permits arbitrary measurable
observations of raw words and exact parameters.

Work on one compact stable continuity cell for the chart inputs,
with frozen atoms having positive actual mass in that cell.  Let
\(\widehat e_m\) be the frozen microscopic exterior in the companion's
resampling calculation, with macroscopic record \(e_m\), and let
\(\mathsf R_m^{\widehat e_m}\) be its independent reference law.
The actual marked conditional law on the cell is exactly
\begin{equation}
 \mathsf Q_m^{\widehat e_m}
   =\frac{W_m^{\widehat e_m}}
          {\mathsf R_m^{\widehat e_m}[W_m^{\widehat e_m}]}
       \mathsf R_m^{\widehat e_m}.
 \label{eq:local-weight}
\end{equation}
This is an identity of raw-word masses, cut and search-mark masses, and
exact retained shifts, established before any observation is selected.
On a sufficiently small such cell, the guards and stable comparison
instructions make every prescribed flexible resolution correct for
every admitted raw configuration.  In the finite bridge, the protected
prefix inventories and exact terminal heights also make the prescribed
suffix resolutions and empty reduction automatic.  Before restriction
to these cells the corresponding indicators remain in the exact weight;
the omitted actual mass is removed only by the exhaustion below.  The
remaining cell tests have continuity boundaries in the macroscopic
variables.  After cancelling factors depending only on the frozen
exterior, the weights are bounded on the cell.  Write \(I_m\) for the
scaled chart input and \(J_m\) for the scaled auxiliary full heights,
noise paths, and whole-contour tests kept in the calculation.  Along
every convergent sequence of good microscopic exteriors whose limiting
frozen value \(e\) belongs to the full-measure continuity set for the
chosen limiting frozen-data law,
their joint reference law satisfies
\begin{equation}
 (I_m,J_m,W_m^{\widehat e_m})
       \Longrightarrow (I,J,W(e,I,J)),
 \qquad
 0<\int W(e,I,J)\,R_e(\dd I,\dd J).
 \label{eq:local-weight-limit}
\end{equation}
Here \(J\) is the limit of \(J_m\).  The limiting weight
depends only on these macroscopic inputs.  These are the
protected-word density and positive-normalizer conclusions of
Section~5 of~\cite{C}, and the resampling and macroscopic-density
conclusions used in the proof of Theorem~7.10 of~\cite{P}.

Adjoining \(Y_m^\alpha\) is a measurable push-forward of both sides
of \eqref{eq:local-weight}, so the finite identity is unchanged.  Let
\(\mathcal P_m\) be the sigma-field of the full permitted protected
package: its raw extended blocks, durations, used-start offsets, exact
forest differences, and instruction.  Its rescaled intrinsic input is
\(I_m\), and the observation \eqref{eq:finite-readout} is
\(\mathcal P_m\)-measurable.  Compactness of \(E_\alpha\) permits an
extraction whose protected-package limit disintegrates as
\(R_e^I(\dd I)K^\alpha(I,\dd y)\).

Let \(D_m\) be the chosen stable feasible cell in the full reference
proposal.  In the cell-restricted identity \eqref{eq:local-weight},
globally
\(W_m^{\widehat e_m}=\ind_{D_m}\overline W_m^{\widehat e_m}\), with
\(\overline W_m^{\widehat e_m}\) extended by zero off the cell and
uniformly bounded after the exterior cancellation.  Thus
\(\ind_{D_m}\mathsf R_m^{\widehat e_m}\) is an unnormalized
subprobability; the package is not conditioned on fitting.  The recorded
lattice and parity conventions put the retained endpoints on the feasible
coset, and the uniform local limit with positive limiting density excludes
dummy endpoints on the retained compact positive ranges for all
sufficiently large \(m\), by \cite[Lemma~5.5]{C} and the proof of
\cite[Proposition~7.9]{P}.  Together with the stable-resolution argument
above, this gives an input-only limiting continuity cell \(D_e\).

Disintegrate \(R_e(\dd I,\dd J)=R_e^I(\dd I)\mathsf J_e(I,\dd J)\),
using the version furnished by the compact increment, fibre-height, and
bridge kernels.  For every bounded continuous \(\psi\), along the good
frozen-exterior sequences above, the source construction implies
\begin{equation}
\begin{gathered}
 \Big\|
 \mathsf R_m^{\widehat e_m}
   [\ind_{D_m}\psi(J_m)\mid\mathcal P_m]
 -\mathsf J_e(\ind_{D_e}\psi)(I_m)
 \Big\|_{L^1(\mathsf R_m^{\widehat e_m})}
 \longrightarrow0,\\
 \mathsf J_e(\ind_{D_e}\psi)(I)
 :=\int \ind_{D_e}(I,j)\psi(j)\mathsf J_e(I,\dd j).
\end{gathered}
\label{eq:reference-conditional-test}
\end{equation}
This is an averaged conditional estimate; rare exact packages may have
large error.  To see the implication from the source, condition first on
the package, fresh ordinary nuisance pieces, and fibre height.  The exact
normalized bridge endpoints differ by \(o_{\mathbb P}(1)\), under the
joint reference law, from those computed from \(I_m\), the scaled fibre
height, and the macroscopic nuisance paths.  Compact-uniform bridge
convergence replaces their bounded continuous conditional tests in
\(L^1\); integration over the nuisance kernels and conditional Jensen
give \eqref{eq:reference-conditional-test}.  These are the compact-cell
ingredients in \cite[Lemma~5.5 and the proof of Lemma~5.7]{C} and the
proof of \cite[Proposition~7.9]{P}.  Here the surrogate bound is used only
for protected extensions and genuinely fresh ordinary nuisance blocks.
For the hard cell indicator one first uses continuous cutoffs.
Approximating from inside and outside, the null limiting boundary makes
the remaining reference \(L^1\) gap vanish, first as \(m\to\infty\) and
then as the cutoff neighborhood shrinks.  The resulting
\(\mathsf J_e(\ind_{D_e}\psi)\) is continuous at \(R_e^I\)-almost every
input.

Multiplying the conditional error by a bounded continuous
\(f(I_m,Y_m^\alpha)\), then using the protected-package extraction,
proves the localized factorization
\begin{equation}
 (I_m,J_m,Y_m^\alpha)_*
   (\ind_{D_m}\mathsf R_m^{\widehat e_m})
   \Longrightarrow
   \ind_{D_e}(I,J)R_e(\dd I,\dd J)\,K^\alpha(I,\dd y).
 \label{eq:reference-disintegration}
\end{equation}
Full heights affect the finite observation only through the exact forest
differences in \(\theta_m^\alpha\), whose scaled values are retained in
\(I_m\); durations remain input variables.  Thus no continuity of the
numerical observation in the contours was used.

Take the next joint limits under these truncated laws.  The corresponding
weight marginal, by \eqref{eq:local-weight-limit} and the same cell
cutoff, places the weight at \(W(e,I,J)\) almost surely.  Since
\(W_m^{\widehat e_m}=\ind_{D_m}W_m^{\widehat e_m}\) and
\(W=\ind_{D_e}W\), the numerator and denominator in
\eqref{eq:local-weight} can be integrated against the truncated laws
without change.
For a bounded continuous test \(\Phi\) of the retained input and output,
boundedness of the weights in \eqref{eq:local-weight} now yields
\begin{equation}
 \lim_m\mathsf Q_m^{\widehat e_m}[\Phi]
 =
 \frac{\displaystyle
   \int \Phi(I,J,y)W(e,I,J)
          R_e(\dd I,\dd J)K^\alpha(I,\dd y)}
      {\displaystyle\int W(e,I,J)R_e(\dd I,\dd J)} .
 \label{eq:weighted-kernel-limit}
\end{equation}
Thus the weight changes the input law and leaves the conditional law
of \(y\) equal to \(K^\alpha(I,\dd y)\).  Including bounded
continuous whole-contour tests in \(J\), followed by a monotone class,
gives the identity conditional on the complete contours and marks.
The conditional-independent attachment of the remaining coordinate
data preserves it.  The later finite-cell exhaustion removes the cell
restriction in the actual law.

For finitely many disjoint buffers, the complete reference packages,
including the added outputs, are independent.  Their joint weak limit
is the product of their joint limits.  The resampling weight for their
union depends only on the combined macroscopic inputs, so
\eqref{eq:weighted-kernel-limit} preserves the product of output
kernels.  Cylinder tests give the countable assertion.  Exact
reconstruction in a common union record makes two descriptions of the
same finite observation equal eventually on activation, and this
equality passes to the joint limit.  For inner and outer descriptions
that are only ordered, the corresponding inequalities pass with the
fixed open losses instead.

For the canonical feasible-order augmentation, fix a compact protection
inside a bounded real buffer, with a fixed larger imaginary buffer.
The traversals of the real buffer meeting the protection form an almost
surely finite family.  Let \(o\) range over its feasible orders, with
positive conditional weights \(w_o(I_{\mathrm{IG}})\).  These weights depend
only on the unclocked imaginary input, not on the real field or the
time unit.  The preceding argument gives the enlarged joint kernel
\(K_o\) at each order.  Keep the actual-order component and sample the
other components independently from their \(K_o\)'s conditional on
the fields.  The vector then has conditional law \(\prod_oK_o\),
with its joint array of labels intact.  Independent auxiliary fills
in disjoint raw-buffer packages preserve their product law.  This is
the feasible-order construction in Section~5 of~\cite{C}, now applied
to the enlarged output.

The compact continuity cells exhaust the actual generic input law.
The companion's reconstruction and guard arguments supply finite
covers, positive gaps, and admissible instructions almost surely.
Finitely many cells and compact parameter truncations omit at most
\(\varepsilon\) of that law.  Their omitted contribution to a bounded
conditional-identity test is at most its bound times \(\varepsilon\).
Letting \(\varepsilon\downarrow0\), and then diagonalizing over
instructions, finite unions and coordinates, proves the simultaneous
identities on activation.  Only input and reconstruction events enter this
exhaustion, so there is no weight depending on a new numerical success
event.

A change of input density \(r(I,J)\) multiplies the limiting
input--output law by that same density.  Bounded approximation and
truncation handle every integrable density.  For changes of units use
\cite[Proposition~5.12]{C} and the closure construction in
\cite[Section~7]{P}.  Expressing the same raw map with time unit \(vm\)
changes the field, time labels, and a labelled frontier length
\(L_\partial\) to
\begin{equation}
 h_v=h-\gamma^{-1}\log v,\qquad t_v=v^{-1}t,\qquad
 L_{\partial,v}=v^{-1/2}L_\partial ,
 \label{eq:indexed-units}
\end{equation}
A raw incidence quantity is unchanged; a normalized output changes
only by the specified new deterministic denominator.  Raw cuts and
incidences remain exact.  Rounding a time unit changes its scaled
labels by \(o(1)\) when that unit diverges.  For varying labels, retain
the normalized reference input law, its positive duration and shift
densities, and all fixed effective units as required by the companion's
indexed construction.  Equation~\eqref{eq:indexed-units} re-expresses
one array; it does not identify differently indexed kernels at an
unchanged field.

For the spatial assertion put
\begin{equation}
 \xi=\bigl(h|_U,\ h^{\mathrm{IG}}|_V
       \pmod{\text{period}},\ \text{area marks in }U,\ O\bigr),
 \label{eq:spatial-input}
\end{equation}
where \(V\) contains a neighborhood of \(\overline U\), and \(O\)
is the finite relative order.  The reconstruction in
\cite[Theorem~7.13]{P}, and its canonical counterpart in
\cite[Section~5]{C}, reads from \(\xi\) the finitely many directed
traversals meeting \(K_*\), their clocked cuts, labelled frontier
relations, forest differences, separating bands, and used
flexible-order instructions.  The containments
\eqref{eq:local-protections} and the intrinsic instruction are
therefore selected from \(\xi\) on activation.  The full preimage
and star conditions make their finite reconstruction exact eventually.

The global labels locating those blocks in \(\mathcal C\) need not be
measurable from \(\xi\).  What is local is the evaluated kernel.
Two realizations of the same intrinsic description use the same
reference blocks, durations, shifts, and observation after relabeling.
Denote that description by \(s(\xi)\), its continuous input by
\(c(\xi)\), and the map on the joint label/output space that sends
time labels to spatial positions, leaving intrinsic numerical values
unchanged, by \(\Phi^\xi_{s(\xi)}\).  The selected spatial kernel is
\begin{equation}
 \mathsf K_\xi
   =\bigl(\Phi^\xi_{s(\xi)}\bigr)_*
       K^{s(\xi)}(c(\xi),\cdot),
 \label{eq:spatial-kernel}
\end{equation}
which is measurable from \(\xi\).  If \(\alpha_*(\mathcal C)\) is the
globally located atlas member selected by this description, then on
its activation event summing the fixed-chart identities over
\(\{\alpha_*=\alpha\}\) and using the tower property gives
\begin{equation}
 \EE[f(Y^{\alpha_*})\mid\mathcal C]=\mathsf K_\xi f,\qquad
 \EE[f(Y^{\alpha_*})\mid\xi]=\mathsf K_\xi f
 \label{eq:spatial-selection}
\end{equation}
for bounded measurable \(f\).  The same computation for products and
common union records proves spatial product and restriction
compatibility.  This selection occurs after the limiting array has
been extracted.
\end{proof}

The observations needed below fit this result directly.  On a finite
retained graph, extremal length with vertex ports and subgraph
confinement is a numerical optimization over a finite edge space.  An
infimum of a quadratic edge energy subject to finitely many vertex-value
bounds is measurable as well: its sublevel tests are finite polynomial
inequalities with real quantifiers.  The same description applies to
a universal inequality over all vertex functions.  A minimum
sufficient dyadic level for such inequalities down to singleton cuts
is one integer-valued function of the same finite record.  These
observations require no continuity in the limiting contours.  Their
spatial meaning always uses the fixed margins and activation above.

\begin{proposition}[Probability-one support on ordinary sphere patches]
\label{prop:sphere-support}
Fix \(T>0\), a joint extraction of ordinary empty-word excursions of
duration \(T\), and matching bilateral or ordinary-field reference
experiments from Proposition~\ref{prop:readout-extension}.  Normalize
a sphere coordinate using the exploration root at infinity and two
independent quantum-area marks at \(0,1\).  Let \(K\Subset U\) be a
protected patch, where \(U\) is a relatively compact disk avoiding
\(0,1\), with a fixed larger imaginary-field buffer.

For each retained coordinate set its limiting output equal to a fixed
cemetery symbol when that coordinate is not activated, and call the
resulting finite or countable vector the activated array.  Let \(B\)
be a measurable event of the local input \eqref{eq:spatial-input} and
this activated array.  If \(B\) has probability zero for the ordinary
reference input with its common conditional observation kernel, it has
probability zero in this sphere extraction.  The assertion retains the
actual additive sphere height and the same deterministic time and
observation units.  It holds simultaneously for a countable protected
atlas and fixed open losses.  Independent additional area-mark pairs
exhaust the auxiliary marks by application to each marked marginal;
the exploration root remains excluded.

Let \(X_m\in[0,\infty]\) be retained through compactification, and
let \(X\in[0,\infty]\) denote its extended-real limit.  If \(X\le C\)
almost surely on activation for a deterministic \(C<\infty\) in every
such ordinary extraction, then, almost surely on activation of that
coordinate in the sphere representation extraction, \(X_m\le C'\)
eventually for every fixed \(C'>C\).  If instead \(X<\infty\) almost
surely on activation in every such ordinary extraction, then, almost
surely on activation of that coordinate in the sphere representation
extraction, \(X_m\) is eventually bounded by a finite constant which may
depend on the sample and the test.
\end{proposition}

\begin{proof}
Let \(\mathsf P_{T,U}\) be the law of the actual duration-\(T\) sphere
real field on \(U\), and fix
\(\mathsf P_{\mathrm{ord},U}:=\operatorname{Law}(G_R|_U)\) with
\(\overline U\Subset B_R\) as above.  The actual-height comparison
\cite[Lemma~5.13 and Proposition~5.12]{C}, or
\cite[Lemma~3.3, as used in Lemma~18.1]{P}, gives
\(\mathsf P_{T,U}\ll\mathsf P_{\mathrm{ord},U}\) for \(T=1\).
For \(T\ne1\) the field is shifted by \(\gamma^{-1}\log T\); a
Cameron--Martin cutoff equal to that constant near \(\overline U\)
gives the same local absolute continuity.

Let \(\mathsf P_{\mathrm{IG}}\) be the law of the entire independent
stationary-phase imaginary field and its unparameterized directed
space-filling curve, and let \(\Pi_h\) be the common conditional law
of the local area marks.  The full imaginary object retains the order
\(O\) in \eqref{eq:spatial-input}, and all the corresponding finite
orders for an array of tests.  For such an array, \(\xi\) and
\(\mathsf K_\xi\) denote its joint local input and kernel.  The input laws satisfy
\begin{equation}
 \mathsf P_{T,U}(\dd h)\,
 \mathsf P_{\mathrm{IG}}(\dd\zeta)\,\Pi_h(\dd\pi)
 \ \ll\
 \mathsf P_{\mathrm{ord},U}(\dd h)\,
 \mathsf P_{\mathrm{IG}}(\dd\zeta)\,\Pi_h(\dd\pi).
 \label{eq:sphere-input-domination}
\end{equation}
The local input \(\xi\) is a measurable function of \((h,\zeta,\pi)\).
Proposition~\ref{prop:readout-extension} gives the same observation
kernel \(\mathsf K_\xi\) on activation on both sides; the cemetery
convention extends it to the activated array everywhere.  If \(r(h)\)
is the input density,
the sphere probability of \(B\) equals
\begin{equation}
 \int r(h)\,\mathsf K_{\xi(h,\zeta,\pi)}(B_{\xi(h,\zeta,\pi)})
       \,\mathsf P_{\mathrm{ord},U}(\dd h)
         \mathsf P_{\mathrm{IG}}(\dd\zeta)\Pi_h(\dd\pi).
 \label{eq:sphere-support-integral}
\end{equation}
The integrand without \(r\) vanishes almost everywhere by the reference
null assumption, so this integral is zero.  A compact patch avoiding
the exploration root has full time preimage in a strict interior of
\((0,T)\), as required for the finite-word kernel.  Countable
intersection gives the simultaneous statement.  Repetition with
fresh area-mark pairs uses each marginal separately, without conditioning
on the other pairs.

Convergence in \([0,\infty]\) to a value at most \(C\) gives an
eventual bound \(C'>C\); convergence to a finite value gives an
eventual finite bound.  Exact reconstruction on activation supplies
the finite-graph interpretation.  This deduction is made in the
representation extraction and supplies no finite-scale rate or
conditional probability.
\end{proof}

For selected neighborhoods of a bilateral word, the needed comparison
is the raw-block subprobability estimate, which can see a discontinuous
lattice observation.

\begin{lemma}[Null tests in protected selected slots]
\label{lem:record-null-transfer}
Take a finite chronologically ordered list of disjoint slots as in
\cite[Proposition~7.16]{P}, each with a left safety band of positive
limiting duration.  Conditional on the incoming past and auxiliary
choices independent of the fresh letters, the entire list of retained
cuts, durations, extensions, and guard/noise anchors is chosen before
fresh letters are inspected.  The auxiliary selections have total
conditional weight at most one.  Retained shifts form a forest in each
coordinate; after the unimodular change to differences from each
component's earliest vertex, at most one new difference per coordinate
is imposed at a slot.  The finite instruction list is genuinely
compatible, and durations, shifts, ordinary increments, and noise
parameters lie in the stated compact fitting ranges.  For the null
transfer below, no slot meets the distinguished center time.  Append observations
\eqref{eq:finite-readout} of those records.

Let \(\mathsf F_m\) be the fresh-word comparison, \(G_m\) the
compatibility event, and \(\mathsf R_m\) the law of the retained input
and observations under its independent raw-block reference experiment.
The observation push-forwards satisfy
\begin{equation}
 (\mathsf F_m|_{G_m})_{\rm obs}\le C\mathsf R_m.
 \label{eq:record-domination}
\end{equation}
Identifying an observation event with its preimage, if a selected law
\(\mathsf S_m\) has bounded relative entropy with
respect to \(\mathsf F_m\), allowing bounded average conditional
entropy with an unchanged history marginal, then
\(\mathsf R_m(B_m)\to0\) implies
\(\mathsf S_m(B_m\cap G_m)\to0\).

In a jointly extracted openly fillable protected chart, every event of
the input and appended observations that is null in the ordinary
completion is null in the selected limit on that chart.  This remains
true for scale-dependent reference laws and a random viewing level
under the joint-extraction hypotheses of
\cite[Proposition~7.18]{P}.
\end{lemma}

\begin{proof}
Proposition~7.16 of~\cite{P} bounds the compatible subprobability law
of the entire raw retained blocks and exact shifts.  A measurable
push-forward preserves that inequality, proving
\eqref{eq:record-domination}.  There is no division by the probability
of compatibility.  Put \(p_m=\mathsf S_m(B_m\cap G_m)\) and
\(q_m=\mathsf F_m(B_m\cap G_m)\).  When \(0<q_m<1\), the binary
relative-entropy inequality gives
\begin{equation}
 p_m\log(1/q_m)
 \le \operatorname{Ent}(\mathsf S_m\mid\mathsf F_m)+\log2 .
 \label{eq:record-entropy}
\end{equation}
If \(q_m=0\), finite entropy already gives \(p_m=0\).  Since
\(q_m\le C\mathsf R_m(B_m)\to0\), the conclusion follows.
Conditional entropy integrates to joint entropy when the history
marginal is unchanged; a tight entropy bound may first be truncated.

For the final assertion, write the reference input and output law as
\(\mu(\dd c)K_c(\dd y)\).  On a compact stable cell the ordinary
completion induces a retained subprobability law
\begin{equation}
 f(c)\,\mu(\dd c)K_c(\dd y),\qquad
 f(c)=\int a(c,v)\ind_{\mathcal T}(c,v)\,\lambda_c(\dd v).
 \label{eq:fillable-density}
\end{equation}
Here \(\lambda_c\) is the conditional law of omitted Brownian or
first-passage filling pieces, \(\mathcal T\) is the strict compatible
filling event, and \(a(c,v)>0\) is the translated endpoint-density
multiplier.  Open fillability gives \(f(c)>0\) on the inputs at issue.
The multiplier depends on the input and filling, so
Proposition~\ref{prop:readout-extension} leaves the enlarged \(K_c\)
unchanged.  If \(B\) is null in this completion, then
\[
 0=\int f(c)K_c(B_c)\,\mu(\dd c)
\]
forces \(K_c(B_c)=0\) for \(\mu\)-almost every openly fillable \(c\).
Restricting first to \(f(c)\ge\varepsilon\) and then letting
\(\varepsilon\downarrow0\) requires no uniform filling probability.
The tight absolute continuity obtained from
\eqref{eq:record-domination}--\eqref{eq:record-entropy} transfers this
null set to the selected limit.  Proposition~7.18 of~\cite{P}
applies the argument in each jointly extracted limit for
scale-dependent reference laws, and hence for mixtures over a viewing
level.  Its off-center and strict-gap hypotheses remain in force.
\end{proof}

We finish with the pointwise consequence of locality used for crossing
directions and for the conductivity density.  Call a jointly
measurable readout \(Y(z)\), \(z\) in an ordinary coordinate domain
\(D\), a \emph{germ} if, for every neighborhood \(V\) of \(z\), it is
a measurable function of the input and observation arrays of activated
strict tests whose buffered supports are compactly inside \(V\).
Descriptions must agree under common refinements.  It is enough that
\(Y(z)\) be a tail limit of such tests as their supports shrink to
\(z\), unchanged by deleting finitely many initial tests.  Examples
include an indicator quantified over all sufficiently small dyadic
disks and a density obtained from local measures on shrinking disks.

\begin{lemma}[Deterministic local germs]
\label{lem:deterministic-germs}
Let \(Y(z)\) be a germ with values in a separable metric space, formed
from the arrays of Proposition~\ref{prop:readout-extension} in an
ordinary field chart.  Assume its description is intrinsic under each
coordinate change under consideration: transporting the fields,
traversals and ports transports \(Y\) by a specified measurable action
on its value space.  Then for planar-area-almost every deterministic
\(z\), \(Y(z)\) is almost surely deterministic.  The deterministic
values obey that action for each fixed affine coordinate change for
which the original and transported ordinary field restrictions are
locally equivalent.  The assertions hold jointly for countably many
germs and coordinate changes.

Suppose in addition that \(D\) is connected and that, for every
\(b\in\mathbb Q^2\) and every pair of disks \(V,V+b\Subset D\), the
descriptions on these disks are related by translation, its action on
the value space is the identity, and common restrictions identify the
descriptions on overlaps.  Then the deterministic value is constant
almost everywhere in \(D\).  In particular, a matrix germ with this
translation property and rotation action \(M\mapsto RMR^{\mathsf T}\)
has a constant matrix \(M_0\) satisfying
\(M_0=RM_0R^{\mathsf T}\) for every permitted rotation \(R\).
\end{lemma}

\begin{proof}
We give the locality argument of \cite[Lemma~8.3]{C} for the enlarged
arrays.  Let \(\mathcal H\) be the sigma field of the entire real and
imaginary fields.  By \cite[Lemma~3.7]{C}, at every deterministic planar
\(z\) the whole-plane exploration has a unique preimage almost surely.
For every larger neighborhood \(V\), all
visits to a sufficiently small neighborhood of \(z\) then lie in one
traversal of \(V\): otherwise properness gives a second limiting
preimage outside that traversal.  Success is readable in a slightly
larger local buffer, since the imaginary field there reconstructs the
finite directed traversal list.  Shrinking the inner neighborhood
exhausts a probability-one event.  On this exhaustion, the small tests
defining a germ need no exterior permutation; their order is the
direction of that traversal.  Proposition~\ref{prop:readout-extension}
therefore implies, for bounded measurable \(\varphi\), that
\begin{equation}
 p_\varphi(z)=\EE[\varphi(Y(z))\mid\mathcal H]
 \label{eq:germ-conditional-mean}
\end{equation}
is measurable from the field restrictions in every neighborhood of
\(z\).  Here the local Poisson auxiliary marks of the kernel construction
are integrated out: their restrictions to disjoint spatial sets are
conditionally independent given \(h\).  This step does not involve the
three canonical area samples used to normalize the sphere.  At distinct
\(z,w\), disjoint sufficiently small real buffers give product
conditional kernels, hence conditional independence of
\(\varphi(Y(z))\) and \(\varphi(Y(w))\) given \(\mathcal H\).

For a bounded real \(\varphi\), conditional Fubini and this product
identity give, for every rational disk \(V\Subset D\),
\begin{equation}
 \EE\!\left[
   \left(\int_V\{\varphi(Y(z))-p_\varphi(z)\}\dd z\right)^2
       \,\middle|\,\mathcal H\right]=0 .
 \label{eq:germ-integral}
\end{equation}
The diagonal in \(V\times V\) has area zero.  A countable disk
exhaustion and Lebesgue differentiation imply
\(\varphi(Y(z))=p_\varphi(z)\) for almost every \(z\), almost surely.
Fubini reverses these quantifiers at almost every deterministic point.

The point germ of the two ordinary fields is trivial.  For a Dirichlet
GFF, the intersection of the Gaussian subspaces generated by
restrictions to shrinking disks around \(z\) is orthogonal to every
smooth function vanishing near \(z\).  Logarithmic cutoffs show that a
point has zero Dirichlet capacity, so these functions are dense and the
intersection is zero.  Decreasing Gaussian projections give triviality
of the completed sigma-field intersection.  The same reasoning applies
to two independent fields; reducing the imaginary field modulo its
period removes information.  Local absolute continuity transfers the
conclusion to the ordinary chart conventions, as in the proof of
Lemma~8.3 of~\cite{C}.  Since \eqref{eq:germ-conditional-mean} is
measurable in every shrinking field restriction, it is deterministic.
A countable family of bounded continuous functions separating points
of the value space proves the claim for \(Y\).

For \(\psi(z)=az+b\), the field transformation is
\begin{equation}
 h\longmapsto h\circ\psi+Q\log|a|,
 \qquad
 h^{\mathrm{IG}}\longmapsto
 h^{\mathrm{IG}}\circ\psi-\chi\arg a,
 \quad
 Q=\frac2\gamma+\frac\gamma2,\quad
 \chi=\frac2\gamma-\frac\gamma2 .
 \label{eq:affine-fields}
\end{equation}
Proposition~5.12 of~\cite{C} identifies the same intrinsic incidence
record in these coordinates and gives local equivalence of the field
laws.  A probability-zero or probability-one statement about \(Y\)
therefore transports by its stated action, so its deterministic value
obeys that action.  Under the additional translation hypothesis, every
bounded continuous function of the deterministic version is invariant
almost everywhere under rational translations whose local buffers fit.
Continuity of translations in local \(L^1\) extends this to all
translations, making each such function constant almost everywhere on
the connected domain.  A countable separating family gives a constant
value of the germ.  The rotation assertion follows with the stated
conjugation action.
\end{proof}
